% ===== PRÉAMBULE CANONIQUE — à coller VERBATIM en tête de CHAQUE .tex =====
\documentclass[11pt,a4paper]{article}
\usepackage[utf8]{inputenc}\usepackage[T1]{fontenc}\usepackage{lmodern}
\usepackage[french]{babel}
\usepackage{amsmath,amssymb,mathtools}\usepackage{icomma}
\usepackage[version=4]{mhchem}          % \ce{...} : équations & formules
\usepackage{siunitx}\sisetup{locale=FR} % \qty{}{}, \si{}, \ang{}
\usepackage{chemfig}                    % \chemfig{...} : structures développées
\usepackage{xcolor}\usepackage{colortbl}
\usepackage{tikz}\usepackage{pgfplots}\pgfplotsset{compat=1.18}
\usepackage[most]{tcolorbox}\usepackage{enumitem}\usepackage{array,booktabs}
\usepackage[a4paper,margin=2.2cm]{geometry}
\usetikzlibrary{arrows.meta,calc,angles,quotes,positioning,patterns,backgrounds,shapes.geometric}
\definecolor{primary}{RGB}{222,49,99}\definecolor{primarydark}{RGB}{150,22,55}
\definecolor{defcol}{RGB}{0,114,178}\definecolor{methcol}{RGB}{52,92,148}
\definecolor{excol}{RGB}{0,135,90}\definecolor{attncol}{RGB}{200,40,55}
\tcbset{boxbase/.style={enhanced,breakable,boxrule=0.9pt,arc=2.5pt,
  left=3mm,right=3mm,top=2.3mm,bottom=2.3mm,fonttitle=\bfseries,coltitle=white}}
% --- boîtes TUTORIEL (noms EXACTS, ne pas renommer) ---
\newtcolorbox{cle}[1][]{boxbase,colback=defcol!6,colframe=defcol,
  title={\textsc{À retenir}\ifx\relax#1\relax\relax\else\textmd{ : \itshape #1}\fi}}
\newtcolorbox{methode}[1][]{boxbase,colback=methcol!7,colframe=methcol,sharp corners,
  title={\textsc{Méthode}\ifx\relax#1\relax\relax\else\textmd{ --- \itshape #1}\fi}}
\newtcolorbox{exemplebox}[1][]{boxbase,colback=excol!5,colframe=excol,
  title={\textsc{Exemple}\ifx\relax#1\relax\relax\else\textmd{ : \itshape #1}\fi}}
\newtcolorbox{piege}[1][]{boxbase,colback=attncol!6,colframe=attncol,
  title={\textsc{Piège}\ifx\relax#1\relax\relax\else\textmd{ : \itshape #1}\fi}}
\newtcolorbox{remarque}[1][]{boxbase,colback=black!4,colframe=black!45,
  title={\textsc{Remarque}\ifx\relax#1\relax\relax\else\textmd{ : \itshape #1}\fi}}
% --- boîtes EXERCICE / CORRIGÉ (noms EXACTS, auto-comptées) ---
\newtcolorbox[auto counter]{exo}[1][]{boxbase,colback=primary!4,colframe=primary,
  title={\textsc{Exercice}~\thetcbcounter\ifx\relax#1\relax\relax\else\textmd{ --- \itshape #1}\fi}}
\newtcolorbox[auto counter]{corrige}[1][]{boxbase,colback=excol!4,colframe=excol,
  title={\textsc{Corrigé}~\thetcbcounter\ifx\relax#1\relax\relax\else\textmd{ --- \itshape #1}\fi}}
\newcommand{\R}{\mathbb{R}}\newcommand{\N}{\mathbb{N}}\newcommand{\Z}{\mathbb{Z}}\newcommand{\ds}{\displaystyle}
% (un \usepackage{fancyhdr}+en-tête est possible ; il est ignoré côté web)
% =========================================================================
\begin{document}
\begin{corrige}[masse moyenne à trois isotopes]
$\overline{A} = \sum_i p_i A_i$.
\begin{enumerate}[label=\alph*)]
  \item $\overline{A} = 0{,}79\times 24 + 0{,}10\times 25 + 0{,}11\times 26 = 18{,}96 + 2{,}5 + 2{,}86
        = 24{,}32$.
  \item $\overline{A} = 0{,}92\times 28 + 0{,}05\times 29 + 0{,}03\times 30 = 25{,}76 + 1{,}45 + 0{,}9
        = 28{,}11$.
\end{enumerate}
\end{corrige}

\begin{corrige}[retrouver une abondance à partir de $\overline{A}$]
On écrit $\overline{A} = pA_1 + (1-p)A_2$ et on résout en $p$.
\begin{enumerate}[label=\alph*)]
  \item $6{,}94 = 6p + 7(1-p) = 7 - p \Rightarrow p = 0{,}06$ : \textbf{$6\,\%$ de $\ce{^{6}Li}$} et
        $94\,\%$ de $\ce{^{7}Li}$.
  \item $63{,}5 = 63p + 65(1-p) = 65 - 2p \Rightarrow 2p = 1{,}5 \Rightarrow p = 0{,}75$ :
        \textbf{$75\,\%$ de $\ce{^{63}Cu}$} et $25\,\%$ de $\ce{^{65}Cu}$.
  \item $10{,}8 = 10p + 11(1-p) = 11 - p \Rightarrow p = 0{,}20$ : \textbf{$20\,\%$ de $\ce{^{10}B}$}
        et $80\,\%$ de $\ce{^{11}B}$.
\end{enumerate}
\end{corrige}

\begin{corrige}[vrai ou faux : isotopes et uma]
\begin{enumerate}[label=\alph*)]
  \item \textbf{\textcolor{excol}{VRAI}} : même $Z$, mais $A$ (donc le nombre de neutrons) diffère.
  \item \textbf{\textcolor{excol}{VRAI}} : c'est exactement la définition de l'uma (= dalton).
  \item \textbf{\textcolor{excol}{VRAI}} : $16$ nucléons, chacun $\approx 1$ uma, donc $\approx 16$ uma.
  \item \textbf{\textcolor{attncol}{FAUX}} : l'uma est \emph{voisine} de la masse d'un proton, mais
        \emph{définie} à partir du carbone 12 ($\tfrac{1}{12}$ de la masse d'un $\ce{^{12}C}$).
  \item \textbf{\textcolor{excol}{VRAI}} : $\SI{1}{uma} = \tfrac{1}{N_{\!A}}\,\si{\gram}
        \approx \SI{1,66054e-24}{\gram}$.
\end{enumerate}
\end{corrige}

\begin{corrige}[avec les masses isotopiques réelles]
Même formule, avec la masse réelle de chaque isotope.
\begin{enumerate}[label=\alph*)]
  \item $\overline{A} = 0{,}758\times 34{,}97 + 0{,}242\times 36{,}97 = 26{,}51 + 8{,}95 = 35{,}45$ uma.
  \item $\overline{A} = 0{,}692\times 62{,}93 + 0{,}308\times 64{,}93 = 43{,}55 + 20{,}00 = 63{,}55$ uma.
\end{enumerate}
On retrouve des valeurs très proches de celles obtenues avec les nombres de masse entiers.
\end{corrige}

\begin{corrige}[deviner l'isotope majoritaire]
La moyenne est plus proche de la masse de l'isotope le plus abondant.
\begin{enumerate}[label=\alph*)]
  \item $\overline{A} = 35{,}5$ est plus proche de $35$ : \textbf{$\ce{^{35}Cl}$} est majoritaire.
  \item $\overline{A} = 63{,}5$ est plus proche de $63$ : \textbf{$\ce{^{63}Cu}$} est majoritaire.
  \item $\overline{A} = 79{,}9$ est (de peu) plus proche de $79$ que de $81$ : \textbf{$\ce{^{79}Br}$}
        est légèrement majoritaire (les deux abondances sont proches).
\end{enumerate}
\end{corrige}

\end{document}
